\documentclass[11pt,letterpaper]{article}
\usepackage[utf8]{inputenc}
\usepackage[T1]{fontenc}
\usepackage{lmodern}
\usepackage[margin=0.75in]{geometry}
\usepackage{booktabs}
\usepackage{xcolor}

\definecolor{azul}{HTML}{172B46}
\definecolor{gris}{HTML}{43556D}
\pagestyle{empty}
\setlength{\parindent}{0pt}
\setlength{\parskip}{6pt}
\renewcommand{\arraystretch}{1.10}

\begin{document}

{\small\bfseries\color{gris} CA-0309 \enspace / \enspace Grupo 5 \enspace / \enspace Ronda 3\par}
{\LARGE\bfseries\color{azul} Garantía y cargo de la cuenta equity-linked\par}
{\small Hombre de 45 años, saldo inicial de 100\,000 u.m. y plazo de cinco años.
1\,000 trayectorias compartidas, semilla 2026 en R y permanencia contractual ponderada.\par}
\vspace{4pt}

\textbf{\color{azul}Comparación por contrato emitido}
{\small\color{gris}\enspace Montos en unidades monetarias.\par}
\smallskip
\begin{tabular}{@{}lrrr@{}}
\toprule
Garantía / cargo & 100\,000 / 2\% & 110\,000 / 2\% & 110\,000 / 3\% \\
\midrule
VP medio de cargos & 9\,194,11 & 9\,194,11 & 13\,524,42 \\
VP medio del complemento & 5\,496,52 & 8\,527,12 & 10\,122,21 \\
NPV medio del proveedor & 3\,697,59 & 667,00 & 3\,402,20 \\
Frecuencia de pérdida & 26,8\% & 36,6\% & 35,6\% \\
CTE95 del proveedor & 24\,480,78 & 31\,413,43 & 29\,976,86 \\
Frecuencia de activación & 38,7\% & 49,1\% & 57,1\% \\
Saldo final medio antes de la garantía & 116\,104,64 & 116\,104,64 & 110\,300,60 \\
\bottomrule
\end{tabular}


\textbf{Elevar la garantía a 110\,000 con cargo del 2\%.}
El VP medio del complemento aumenta 3\,030,59 u.m. y el NPV medio disminuye por el mismo importe.
La frecuencia de pérdida sube 9,8 puntos porcentuales y el CTE95 aumenta 6\,932,66 u.m.

\textbf{Subir el cargo al 3\% con garantía de 110\,000.}
El proveedor recibe 4\,330,30 u.m. más en VP medio de cargos y asume 1\,595,10 u.m. más en VP medio de complemento.
Mejoran el NPV medio, la frecuencia de pérdida y el CTE95. El saldo final medio antes de la garantía baja 5\,804,04 u.m. y la frecuencia de activación sube 8 puntos porcentuales.

\vspace{3pt}
\textbf{\color{azul}Trayectoria 1}
{\small\color{gris}\enspace Cierre del año 5 y resultado del proveedor.\par}
\smallskip
\begin{tabular}{@{}lrrr@{}}
\toprule
Garantía / cargo & 100\,000 / 2\% & 110\,000 / 2\% & 110\,000 / 3\% \\
\midrule
Saldo al cierre del año 5 & 92\,616,64 & 92\,616,64 & 87\,986,76 \\
Complemento al vencimiento & 7\,383,36 & 17\,383,36 & 22\,013,24 \\
VP de cargos & 8\,801,98 & 8\,801,98 & 12\,958,18 \\
VP del complemento & 5\,118,63 & 12\,051,29 & 15\,261,02 \\
NPV del proveedor & 3\,683,35 & -3\,249,31 & -2\,302,84 \\
\bottomrule
\end{tabular}


El primer rendimiento fue de 12,5344\%. Con cargo del 2\% se invirtieron 98\,000 u.m. y el saldo cerró en 110\,283,67.
Con cargo del 3\% se invirtieron 97\,000 u.m. y el saldo cerró en 109\,158,32.
En los tres diseños se activa la garantía. En el primero, el VP de los cargos supera el VP del complemento y el proveedor obtiene un NPV positivo.

\vspace{3pt}
\textbf{\color{azul}Recomendación al Comité.}
Elegir garantía de 100\,000 u.m. y cargo del 2\% por su mayor NPV medio y menor pérdida de cola entre los diseños evaluados.
Condicionar la oferta a contar con recursos líquidos para complementos en escenarios adversos y a un límite de pérdida de cola aprobado por el Comité.
Dar seguimiento al CTE95 por contrato, que promedia las 50 mayores pérdidas.
La estimación se limita a una cuenta representativa y a retornos anuales independientes con parámetros fijos.

\vfill
{\footnotesize\color{gris} Fuente: guía revisada del grupo 5, secciones E, 3 y 4, y tabla de mortalidad suministrada de Costa Rica 2023.\par}

\end{document}
